\section{Simulation Results and Benchmarking}
\label{sec:results}

\begin{table*}[t]
  \centering
  \caption{Core disturbance-family comparison for the benchmark study.
    RMSE values are reported in mHz and computed as family means of
    scenario means. The uncertainty term is the mean within-scenario Monte
    Carlo standard deviation for the winning estimator. $\Delta$ is the RMSE gap
    between the winner and the runner-up.}
  \label{tab:family_summary}
  \scriptsize
  \setlength{\tabcolsep}{2.8pt}
  \renewcommand{\arraystretch}{1.1}
  \begin{tabular}{l c l c l c l}
    \toprule
    Family & Count & RMSE winner & RMSE [mHz] & Runner-up & $\Delta$ [mHz] & Lowest $T_{\mathrm{trip}}$ [s] \\
    \midrule
    IEEE ramps & 9 & SOGI-PLL & $19.4 \pm 6.4$ & Koopman & 2.7 & 0 (multiple) \\
    Mag.\ steps & 7 & EKF & $0.388 \pm 0.143$ & RA-EKF & 1.45 & 0 (multiple) \\
    Modulation & 3 & RA-EKF & $17.8 \pm 21.4$ & Koopman & 13.7 & 0 (multiple) \\
    Phase jumps & 3 & SOGI-FLL & $174 \pm 67.5$ & PLL & 4.03 & ZCD: 0.0194 \\
    OOB interf. & 1 & EKF & $13.6 \pm 11.1$ & RA-EKF & 7.35 & EKF, RA-EKF: 0 \\
    IBR harmonics & 3 & EKF & $70.5 \pm 27.3$ & TFT & 4.13 & TFT, Prony: 0.00947 \\
    IBR ringdown & 5 & UKF & $313 \pm 23.2$ & EKF & 6.67 & ESPRIT: 0.0989 \\
    IBR multi-event & 1 & PLL & $173.5 \pm 66.7$ & UKF & 2.74 & IpDFT: 0.0439 \\
    \bottomrule
  \end{tabular}
\end{table*}

\subsection{Family-Level Accuracy and Trip-Risk}

Table~\ref{tab:family_summary} summarizes family-level RMSE and trip-risk
results. The ranking changes across disturbance families. SOGI-PLL gives the
lowest mean RMSE on IEEE ramp cases (19.4\,mHz). EKF leads magnitude steps
(0.388\,mHz) and out-of-band interference (13.6\,mHz), and RA-EKF leads the
modulation family (17.8\,mHz). Zero-trip-risk ties are common in ramps,
magnitude steps, and modulation, so those families are separated mainly by
accuracy rather than by protection-oriented risk.

Across the 34 benchmark scenarios, nine different estimators win at least one
RMSE scenario, and no estimator wins more than 10 of them. The
magnitude-step family is the clearest exception: EKF wins all seven
magnitude-step scenarios and the family mean as well. This is also the family
with the clearest aggregate margin: EKF leads RA-EKF by 1.45\,mHz, while the
winner's mean within-scenario spread is 0.143\,mHz.

Most other families are tighter. In ramps, modulation, phase jumps, harmonics,
ringdown, and out-of-band interference, the family-level margin is smaller than
or comparable to the same spread. For those families, the winner should be
interpreted as a near-front method rather than as a decisive leader.

Trip-risk is less discriminating in the clean standard cases. In 23 of the 34
scenarios, the minimum trip-risk is zero, often shared by 10 to 13 estimators.
Phase jumps are the main exception. SOGI-FLL gives the lowest mean RMSE across
the phase-jump family (174\,mHz), but ZCD gives the lowest mean trip-risk
($0.0194$\,s). Even inside the standard test set, RMSE and trip-risk are not
interchangeable objectives.

\subsection{Composite IBR Stress Cases}

The largest ranking changes appear under composite stress. The hardest
benchmark scenarios, measured by median RMSE across estimators, are the IBR
multi-event sequence (826\,mHz), the two $60^\circ$ phase-jump cases
(675\,mHz and 637\,mHz), and the noisy IBR ringdown variants
(474--528\,mHz). Median RMSE in these cases exceeds the magnitude-step family
by roughly three orders of magnitude.

The IBR multi-event sequence is the clearest example. PLL gives the lowest
RMSE at 173.5\,mHz. UKF trails by only 2.74\,mHz on the family mean. EKF
remains close at 176.4\,mHz while reducing trip-risk to $0.0874$\,s. IpDFT
gives the lowest trip-risk ($0.0439$\,s) but at the cost of a much larger peak
error (2550\,mHz). The ranking therefore depends on the deployment objective:
minimizing average error, minimizing time above the trip threshold, or bounding
transient severity.

The same metric split appears in the other IBR-specific families. EKF gives the
lowest mean RMSE across the harmonic family (70.5\,mHz), but TFT and Prony
give the lowest mean trip-risk there ($0.00947$\,s). In the ringdown family,
UKF gives the lowest mean RMSE (313\,mHz), whereas ESPRIT gives the lowest
mean trip-risk ($0.0989$\,s) at a much higher computational cost.

Among the 12 scenarios with a unique trip-risk winner, nine differ from the
RMSE leader. Those disagreements concentrate in the $60^\circ$ phase-jump
cases and in the harmonic, ringdown, and multi-event IBR cases. This metric
split is why a single scalar ranking is insufficient for
protection-oriented selection.

\subsection{Global Aggregation and Metric Dependence}

Family-level reading remains the primary analysis, but global summaries still
help as secondary evidence. Their interpretation depends on the aggregation
rule. Mean-rank summaries across the 34 benchmark scenarios favor TFT on RMSE
($4.38$) and peak error ($4.32$), with RA-EKF close behind on both ($4.62$
and $4.59$). Simple averaging over the 34 scenarios instead favors RA-EKF and
EKF, with mean RMSE values of 116\,mHz and 123\,mHz, followed by UKF at
131\,mHz. The difference is informative: mean rank rewards repeated near-front
finishes, whereas raw averaging is driven more strongly by the hardest
families.

Consistency statistics make that distinction concrete. EKF is the most frequent
accuracy leader, with 10 outright RMSE wins, 19 RMSE top-three finishes, and
20 peak-error top-three finishes across the 34 scenarios. TFT wins only two
RMSE scenarios, yet it reaches the RMSE top three in 16 scenarios and the
peak-error top three in 17. Trip-risk favors a different group again: IpDFT
and ESPRIT enter the trip-risk top three in 25 and 24 scenarios, respectively.
Global summaries are therefore screening tools rather than replacements for
family-level evidence.

\begin{table}[t]
  \centering
  \caption{Compact global screening indicators across the 34 benchmark
    scenarios. These statistics summarize different aggregation rules and are
    intended as screening aids rather than as a universal ranking.}
  \label{tab:global_screening}
  \scriptsize
  \setlength{\tabcolsep}{3pt}
  \renewcommand{\arraystretch}{1.08}
  \begin{tabular}{L{1.85cm} L{2.1cm} c}
    \toprule
    Global view & Best estimator & Value \\
    \midrule
    Mean RMSE & RA-EKF & 116\,mHz \\
    Mean-rank RMSE & TFT & 4.38 \\
    RMSE wins & EKF & 10 / 34 \\
    RMSE top-three & EKF & 19 / 34 \\
    Peak-error top-three & EKF & 20 / 34 \\
    Trip-risk top-three & IpDFT & 25 / 34 \\
    \bottomrule
  \end{tabular}
\end{table}

Table~\ref{tab:global_screening} makes the aggregation split explicit. The
rank-based view favors repeated near-front finishes, the mean-RMSE view is more
sensitive to the hardest families, and the hit-count view rewards broad
coverage. No estimator leads all three views at once. Across the 544
estimator--scenario summaries, RMSE correlates strongly with trip-risk
($\rho=0.890$) and settling time ($\rho=0.867$), but only weakly and
negatively with CPU cost ($\rho=-0.202$).

\subsection{Computational Cost}

CPU cost changes the practical ranking again. Figure~\ref{fig:dashboard_mega}
gathers representative time-domain cases together with the global summaries.
EKF and RA-EKF occupy the low-error middle-cost region with mean CPU costs of
$0.813\,\mu$s and $1.61\,\mu$s per sample. SOGI-PLL and SOGI-FLL remain much
cheaper ($0.0867\,\mu$s and $0.0348\,\mu$s) and therefore stay relevant for
embedded protection even when they lose some accuracy. Window-based methods
such as TFT and IpDFT improve trip-risk in some composite cases at about
$3$--$4\,\mu$s per sample. Koopman and ESPRIT occupy a much higher-cost
region, with mean costs of $390\,\mu$s and $2838\,\mu$s, without a matching
global RMSE gain. On the RMSE--CPU plane, EKF, RA-EKF, SOGI-FLL, SOGI-PLL, and
ZCD are Pareto-optimal; in practice, the protection-relevant subset is
\{SOGI-FLL, SOGI-PLL, EKF, RA-EKF\}, since ZCD reaches the front only by
near-zero cost and not by usable composite-stress accuracy.

For deployment screening within this benchmark, estimator choice remains tied
to disturbance family, operating objective, and compute budget. SOGI-PLL is
strongest on ramps, EKF on magnitude steps and out-of-band interference,
RA-EKF on modulation, and PLL on the composite IBR event. Figure~\ref{fig:dashboard_mega}(g)
keeps SOGI-FLL, SOGI-PLL, EKF, and RA-EKF in the main low-cost cluster below
$2\,\mu$s per sample, whereas TFT and IpDFT move rightward in cost but improve
trip-risk behavior in the composite cases. Those are conditional results, not a
universal hierarchy.

\begin{figure*}[!t]
  \centering
  \includegraphics[width=\textwidth,height=0.80\textheight,keepaspectratio]{Fig2_Mega_Dashboard.pdf}
  \caption{Representative benchmark views. Panels (a)--(d) show single-scenario
    responses summarized over the same 60 Monte Carlo runs used throughout the
    study. Panel (e) gives a disturbance scorecard across benchmark families.
    Panel (f) shows the scenario-level RMSE distribution grouped by estimator
    family. Panels (g) and (h) show global cost-versus-accuracy and
    cost-versus-trip-risk views; the dashed vertical line in both panels marks
    the 100-$\mu$s/sample real-time budget implied by the 10-kHz acquisition
    rate. The dashboard shows that estimator selection depends on disturbance
    family, risk metric, and compute budget.}
  \label{fig:dashboard_mega}
\end{figure*}
